\section{Conclusion}

This work treated agentic financial QA as a resource-allocation hypothesis and
allowed a prespecified development gate to reject the current implementation.
CoT/PoT formed the strongest observed development quality--cost point, while
Selective Nexus was less accurate and more expensive; the fallback escalated
70\% of the objective set yet underperformed the static path. We therefore did
not open the disjoint 900-example final manifest. The broader lesson is not that
iterative agents always fail, but that routing cannot manufacture value from an
expert with little complementary accuracy. Establishing that complementarity
should precede difficulty modeling, threshold tuning, and confirmatory testing.
